\documentclass{nfr}

\title{Data Management Plan}
\author{Zhongli Wang}

\begin{document}
\maketitle

\section*{1. Types of Data}

The project will generate three main types of data:
\begin{itemize}
    \item \textbf{Simulation data}: outputs from agent-based and network-diffusion models (aggregated indicators, time-series, adoption distributions).
    \item \textbf{Text-based datasets}: publicly available policy documents, media articles, and institutional reports processed using LLM-assisted workflows.
    \item \textbf{Qualitative inputs}: notes from expert consultations or stakeholder workshops (non-personal and anonymised).
\end{itemize}
No sensitive personal data will be collected; all data involve aggregate or publicly available sources.

\section*{2. Documentation and Metadata}

All datasets will include clear documentation describing: (i) purpose, generation method, and processing steps; (ii) model parameters and version control information; (iii) metadata using Dublin Core where applicable.
Code and workflows will be documented using README files and reproducible scripts.

\section*{3. Storage and Security}

Data will be stored on secure University of Turku servers with regular backups.  
Simulation code will be maintained in a private Git repository with access control.  
Secondment-related files at Lucidminds AI will follow the same security standards and EU data protection guidelines.  
No personal or confidential data will be stored or transferred.
All research data and materials necessary for reproducibility will be preserved for at least ten years on institutional repositories. Long-term maintainability of code and simulation models will be ensured through proper versioning and documentation.


\section*{4. Ethical and Legal Compliance}

The project complies with GDPR, University of Turku ethical guidelines, and Horizon Europe open science requirements.  
Only publicly available or fully anonymised data will be used.  
Textual data derived from policy documents or news sources will follow fair-use and licensing conditions.

\section*{5. Sharing, Reuse, and Open Access}

Where legally permissible:
\begin{itemize}
    \item processed datasets (not raw corpora), metadata, and simulation outputs will be deposited in a trusted repository (Zenodo or UTU Research Data Portal);
    \item code for models and analysis workflows will be released under an open-source licence (e.g., MIT or Apache 2.0);
    \item publications will follow open-access guidelines.
\end{itemize}
Restrictions apply only when data licenses prevent redistribution (e.g., proprietary news corpora).

% \section*{6. Preservation}


\end{document}
